\documentclass[11pt,a4paper]{article}

% -------------------------------------------------
% Sprache / Kodierung
% -------------------------------------------------
\usepackage[T1]{fontenc}
\usepackage[utf8]{inputenc}
\usepackage[ngerman]{babel}
\usepackage{lmodern}
\usepackage{microtype}

% -------------------------------------------------
% Mathematik
% -------------------------------------------------
\usepackage{amsmath,amssymb,amsthm,mathtools}
\usepackage{mathrsfs}
\usepackage{bm}

% -------------------------------------------------
% Layout / Grafik
% -------------------------------------------------
\usepackage[a4paper,margin=2.6cm]{geometry}
\usepackage{graphicx}
\usepackage{xcolor}
\usepackage{hyperref}
\usepackage{enumitem}
\usepackage{booktabs}
\usepackage{tikz}
\usetikzlibrary{arrows.meta,calc,positioning}

% -------------------------------------------------
% Theorem-Umgebungen
% -------------------------------------------------
\newtheorem{definition}{Definition}[section]
\newtheorem{satz}[definition]{Satz}
\newtheorem{lemma}[definition]{Lemma}
\newtheorem{proposition}[definition]{Proposition}
\theoremstyle{remark}
\newtheorem{bemerkung}[definition]{Bemerkung}
\newtheorem{beispiel}[definition]{Beispiel}
\newtheorem{hypothese}[definition]{Hypothese}

% -------------------------------------------------
% Makros
% -------------------------------------------------
\newcommand{\N}{\mathbb{N}}
\newcommand{\Z}{\mathbb{Z}}
\newcommand{\R}{\mathbb{R}}
\newcommand{\C}{\mathbb{C}}
\newcommand{\T}{\mathbb{T}}
\newcommand{\OO}{\mathbb{O}}
\newcommand{\Acal}{\mathcal{A}}
\newcommand{\Hcal}{\mathcal{H}}
\newcommand{\Dcal}{\mathcal{D}}
\newcommand{\Fcal}{\mathcal{F}}
\newcommand{\tr}{\operatorname{tr}}
\newcommand{\diag}{\operatorname{diag}}
\newcommand{\Hol}{\operatorname{Hol}}
\newcommand{\sgn}{\operatorname{sgn}}

% -------------------------------------------------
% Titel
% -------------------------------------------------
\title{Primzahruhr, Aharonov--Bohm-Holonomie, Feinstruktur-Shift und Chern-Phasen\\
	\large Ein diskretes Gauge-Dirac-Modell im Bamberg-Modell}

\author{Thomas Hoffbauer\\Bamberg, Deutschland}
\date{\today}

\begin{document}
	\maketitle
	
	\begin{abstract}
		Wir entwickeln ein eigenständiges diskretes Modell, in dem Tetraederschichten, eine gewichtete Primzahruhr, ein gauge-kovariante Dirac-Operator auf einem Kanalraum, echte Aharonov--Bohm-Flüsse auf elementaren Schleifen und ein Feinstruktur-Shift der Form $2\pi/137$ zu einem gemeinsamen Rahmen verbunden werden. Der Trägerraum entsteht aus den beobachteten Zuständen $(a,b,c)$ eines diskreten Kanalgitters; seine Eulerzahl beträgt in der berücksichtigten 3-dimensionalen Zellstruktur $\chi=1$. Auf diesem Träger wird ein anisotroper Gauge-Dirac-Operator mit E0/ABC-Split, Rough-Term und primzahluhrbasiertem Kantenfilter formuliert. Echte Plaquettenflüsse verbessern die lokale spektrale Feinstruktur der ersten 1000 Riemann-Nullstellen deutlich. Besonders bemerkenswert ist, dass ein Feinstruktur-Shift des Plaquettenflusses von der Form $\Phi\mapsto \Phi+2\pi/137$ die beste gefundene Spacing-Korrelation numerisch signifikant erhöht. In einer periodisierten 2-Band-Version desselben Grundgedankens treten Berry-Krümmung und nichttriviale Chern-Phasen auf. Damit ergibt sich eine geschlossene Kette aus diskreter Geometrie, Schleifenholonomie, spektraler Feinstruktur und topologischer Bandantwort.
	\end{abstract}
	
	\section{Einleitung}
	
	Die leitende Idee dieses Artikels ist, dass mehrere zuvor getrennt behandelte Motive --- Tetraederschichten, Primzahlfamilien, diskrete Holonomien, Gauge-Dirac-Operatoren, Feinstrukturkalibrierung und Chern-Zahlen --- zu einem gemeinsamen mathematisch-physikalischen Rahmen zusammengeführt werden können. Der Ausgangspunkt ist eine diskrete Schalenstruktur, in der natürliche Zahlen und daraus abgeleitete Spektralgrößen auf Tetraederschichten organisiert werden. Diese Schichten tragen eine gewichtete Primzahruhr, die in den aktiven Familien $A,B,C$ eine kanalartige Uhrphase und -amplitude definiert.
	
	Der nächste Schritt besteht darin, diese Uhr nicht als bloße Zähluhr, sondern als \emph{Kantenfilter} in einen diskreten Gauge-Dirac-Operator einzubauen. Dadurch entsteht ein numerisch wirksames Modell, in dem nicht diagonale Potentiale, sondern Kantenphasen, Kantenfilter und echte Schleifenholonomien die wesentliche Feinstruktur tragen. Echte Aharonov--Bohm-Flüsse auf elementaren Plaquetten des Kanalgitters verbessern die lokale spektrale Approximation der ersten 1000 Riemann-Nullstellen. Eine zusätzliche Feinstrukturkalibrierung der Form $2\pi/137$ wirkt in diesem Zusammenhang nicht als bloße Kanalphase, wohl aber als mikroskopische Verschiebung des Plaquettenflusses.
	
	In einer periodisierten Minimalversion desselben Grundgedankens werden schließlich Berry-Krümmung und Chern-Zahl sichtbar. Dadurch treten Eulerzahl des Trägers, Aharonov--Bohm-Holonomie, Feinstrukturverschiebung und von-Klitzing-artige Topologie erstmals in einem gemeinsamen Modell auf.
	
	\section{Tetraederschichten und Primzahruhr}
	
	\subsection{Tetraederschichten}
	
	Wir beginnen mit den Tetraederzahlen
	\begin{equation}
		T_L = \frac{L(L+1)(L+2)}{6},\qquad L\in\N.
	\end{equation}
	
	Die $L$-te äußere Schale ist das Intervall
	\begin{equation}
		(T_{L-1},T_L].
	\end{equation}
	
	Diese Schichten dienen als diskrete Hüllen, auf denen glatte Zahlen, Defektzahlen und Primzahlen beobachtet werden. Die Parität von $L$ wird heuristisch als Phasenindikator gelesen:
	\begin{itemize}
		\item ungerade Schichten: Tri-Effekt,
		\item gerade Schichten: oktonischer Defekt.
	\end{itemize}
	
	\subsection{Gewichtete Primzahruhr}
	
	Für jede Schicht betrachten wir die ungeraden Primzahlen und zerlegen sie nach ihren Restklassen modulo $12$ in die vier Familien
	\begin{equation}
		E\equiv 1,\qquad A\equiv 5,\qquad B\equiv 7,\qquad C\equiv 11\pmod{12}.
	\end{equation}
	
	Anstelle bloßer Anzahlen verwenden wir eine gewichtete Uhr über die Summen
	\begin{equation}
		W_E = \sum_{p\in E\cap \text{Schicht}} \log p,
		\quad
		W_A = \sum_{p\in A\cap \text{Schicht}} \log p,
		\quad
		W_B = \sum_{p\in B\cap \text{Schicht}} \log p,
		\quad
		W_C = \sum_{p\in C\cap \text{Schicht}} \log p.
	\end{equation}
	
	Da $E$ im Folgenden als Hintergrundkanal verstanden wird, definieren wir die aktive Uhr auf dem Dreiersystem $A,B,C$. Dazu setzen wir
	\begin{equation}
		\zeta_{\mathrm{clock}} = \frac{W_A e^{i\theta_A}+W_B e^{i\theta_B}+W_C e^{i\theta_C}}{W_A+W_B+W_C},
	\end{equation}
	mit
	\begin{equation}
		\theta_A=0,
		\qquad
		\theta_B=\frac{2\pi}{3},
		\qquad
		\theta_C=\frac{4\pi}{3}.
	\end{equation}
	
	\begin{definition}[Primzahruhr]
		Die \emph{Primzahruhr} einer Schicht ist die komplexe Größe $\zeta_{\mathrm{clock}}$. Ihre Amplitude
		\begin{equation}
			u := |\zeta_{\mathrm{clock}}|
		\end{equation}
		heißt Uhramplitude, ihr Argument die Uhrphase.
	\end{definition}
	
	\begin{bemerkung}
		Numerisch erwies sich nicht die dominante Familie allein, sondern insbesondere die Uhramplitude als relevante Größe für feinere Spektraltests.
	\end{bemerkung}
	
	\section{Diskreter Kanalraum und E0/ABC-Split}
	
	Wir betrachten Zustände
	\begin{equation}
		x=(a,b,c)\in\N_0^3
	\end{equation}
	als diskrete Kanalzustände. Zu jedem Zustand definieren wir den E0/ABC-Split durch
	\begin{equation}
		e_0(x):=\min(a,b,c),
	\end{equation}
	und
	\begin{equation}
		\alpha=a-e_0,
		\qquad
		\beta=b-e_0,
		\qquad
		\gamma=c-e_0.
	\end{equation}
	
	Damit wird
	\begin{equation}
		(a,b,c)=e_0(1,1,1)+(\alpha,\beta,\gamma)
	\end{equation}
	in einen isotropen Kern und eine anisotrope Kanalexcitation zerlegt.
	
	Zusätzlich tragen die Zustände eine Rough-Komponente $\log m(x)$, die aus der Projektion von Spektralwerten auf glatte Träger stammt.
	
	\section{Diskreter Träger und Eulerzahl}
	
	Der aktuelle Trägerraum ist ein endlicher diskreter Zellkomplex auf dem beobachteten Kanalgitter. Die Knoten sind die tatsächlich vorkommenden Zustände $x=(a,b,c)$. Kanten entstehen zwischen Nachbarzuständen mit Gitterabstand $1$ in einer Koordinate. Plaquetten sind vollständige Einheitsquadrate in den Koordinatenebenen, 3-Zellen vollständige Einheitswürfel.
	
	Für den konkret untersuchten Komplex ergaben sich numerisch:
	\begin{equation}
		V=83,
		\qquad
		E=153,
		\qquad
		F=80,
		\qquad
		C=9.
	\end{equation}
	
	Daraus folgen zwei natürliche Eulerzahlen:
	\begin{equation}
		\chi_{2d}=V-E+F=10,
	\end{equation}
	und für den 3-dimensionalen Zellkomplex
	\begin{equation}
		\chi_{3d}=V-E+F-C=1.
	\end{equation}
	
	\begin{bemerkung}
		Die Zahl $\chi_{3d}=1$ spricht dafür, dass der aktuelle endliche Träger topologisch einem zusammenhängenden Kernkörper nähersteht als einem torischen Objekt. Für periodische Bandtopologie muss daher eine periodisierte Version betrachtet werden.
	\end{bemerkung}
	
	\section{Gauge-Dirac-Operator mit Primzahruhr-Kantenfilter}
	
	\subsection{Kantenamplituden und Phasen}
	
	Der diskrete Gauge-Operator wird auf den drei elementaren Richtungen
	\begin{equation}
		e_A=(1,0,0),\qquad e_B=(0,1,0),\qquad e_C=(0,0,1)
	\end{equation}
	definiert. Eine orientierte Kante $x\to x+e_j$ erhält ein komplexes Hopping
	\begin{equation}
		w_j(x)=t_j(x)e^{i\phi_j(x)}.
	\end{equation}
	
	Der Betrag ist von der Form
	\begin{equation}
		t_j(x)=t_0\bigl(1+\eta_0 e_0 + \eta_1\chi_j + \eta_2\log m\bigr)
		\,F(x,x+e_j),
	\end{equation}
	mit
	\begin{equation}
		\chi_A=\alpha,
		\qquad
		\chi_B=\beta,
		\qquad
		\chi_C=\gamma,
	\end{equation}
	und einem Uhrfilter
	\begin{equation}
		F(x,y)=\exp\!\bigl(-\kappa(\bar\nu(x,y)-\nu_*)\bigr),
		\qquad
		\bar\nu(x,y)=\frac{\nu(x)+\nu(y)}{2}.
	\end{equation}
	
	\begin{bemerkung}
		Numerisch erwies sich die Primzahruhr-Amplitude nicht als guter diagonaler Energieterm, wohl aber als Kantenfilter, also als Übergangsmodulation.
	\end{bemerkung}
	
	Die Phasen enthalten ABC- und Rough-Anteile. In der hier gewählten Minimalform werden drei Richtungsphasen verwendet,
	\begin{equation}
		\phi_A \sim \mu_A\alpha + \rho_A \log m,
		\qquad
		\phi_B \sim \mu_B\beta + \rho_B \log m,
		\qquad
		\phi_C \sim \mu_C\gamma + \rho_C \log m.
	\end{equation}
	
	\subsection{Dirac-Form}
	
	Mit gauge-kovarianten Differenzen $\nabla_j^{(\Acal)}$ definieren wir den diskreten Dirac-Operator
	\begin{equation}
		\Dcal
		=
		i\sum_{j\in\{A,B,C\}} \sigma_j\nabla_j^{(\Acal)} + M,
	\end{equation}
	mit einem milden Massenterm
	\begin{equation}
		M(x)=m_0+m_E e_0(x)+m_{\mathrm{aniso}}\,\mathrm{Aniso}(x)+m_r\log m(x).
	\end{equation}
	
	Dabei bezeichnet $\sigma_j$ die Pauli-Matrizen. Der wesentliche numerische Befund lautet: Nicht diagonale Potentiale, sondern Kantenphasen und Kantenfilter tragen die wesentliche Feinstruktur.
	
	\section{Aharonov--Bohm-Fluss auf echten Plaquetten}
	
	\subsection{Trivialer und echter Fluss}
	
	Ein globaler Richtungsphasen-Shift genügt nicht, um einen echten Aharonov--Bohm-Effekt zu erzeugen. Erst wenn ein Fluss auf elementaren Schleifen nicht weggeeicht werden kann, entsteht eine echte Holonomie. In einem Landau-artigen Gittergauge wird der Fluss in der $AB$-Ebene durch eine Zusatzphase auf den $B$-Kanten realisiert,
	\begin{equation}
		\phi_B^{(\mathrm{AB})}(x)=\Phi\,a,
	\end{equation}
	so dass der Umlauf um elementare Plaquetten einen effektiven Fluss trägt.
	
	\begin{definition}[Plaquettenholonomie]
		Für eine elementare Schleife $P$ sei
		\begin{equation}
			\Hol(P)=\exp\Bigl(i\sum_{(x,j)\in P}\phi_j(x)\Bigr).
		\end{equation}
		Die zugehörige Flussphase ist die Argumentsumme im Exponenten.
	\end{definition}
	
	\subsection{Numerischer Effekt}
	
	Ein echter Plaquettenfluss verbessert die Spacing-Struktur numerisch gegenüber dem Fall ohne Schleifenfluss. Besonders wichtig ist, dass der optimale Fluss nicht bei $\Phi=0$ liegt. Im untersuchten Modell ergab sich ein deutlicher lokaler Feinstrukturgewinn bei bestimmten Flusswerten.
	
	\section{Feinstruktur-Shift $2\pi/137$}
	
	\subsection{Warum ein bloßer Kanaloffset nicht genügt}
	
	Eine konstante Zusatzphase auf einem einzelnen Kanal erwies sich numerisch als zu harmlos; sie ist in der einfachsten Form spektral nahezu wirkungslos. Dagegen wurde eine Feinverschiebung des \emph{echten Plaquettenflusses}
	\begin{equation}
		\Phi_{\mathrm{eff}} = \Phi + \frac{2\pi}{137}
	\end{equation}
	produktiver.
	
	\begin{bemerkung}
		Hier wird $1/137$ nicht als rohe Geschwindigkeit und auch nicht bloß als Kanalphase interpretiert, sondern als feine Justierung der Holonomie selbst.
	\end{bemerkung}
	
	\subsection{Numerischer Haupteffekt}
	
	Mit dieser Verschiebung stieg die beste gefundene Spacing-Korrelation im endlichen Modell erheblich an. Im untersuchten Datensatz ergab sich:
	\begin{equation}
		\text{ohne Shift:}
		\quad
		\mathrm{corr}_{\mathrm{spacing}}\approx 0.322,
	\end{equation}
	gegenüber
	\begin{equation}
		\text{mit Shift }\Phi\mapsto \Phi+2\pi/137:
		\quad
		\mathrm{corr}_{\mathrm{spacing}}\approx 0.562.
	\end{equation}
	
	Der globale affine RMSE blieb dabei ungefähr auf demselben Niveau. Das spricht dafür, dass der Feinstruktur-Shift nicht das grobe Wachstum, wohl aber die lokale Interferenzfeinstruktur verbessert.
	
	\section{Periodisches Minimalmodell und Chern-Zahl}
	
	\subsection{Bloch-Hamiltonian}
	
	Um die globale topologische Antwort zu erfassen, betrachten wir ein periodisches 2-Band-Modell auf dem Torus $\T^2$. Der Bloch-Hamiltonian sei
	\begin{equation}
		H(k)=d_1(k)\sigma_1+d_2(k)\sigma_2+d_3(k)\sigma_3,
	\end{equation}
	mit
	\begin{equation}
		d_1(k)=v_A\sin k_x,
		\qquad
		d_2(k)=v_B\sin k_y,
	\end{equation}
	und
	\begin{equation}
		d_3(k)=M+v_C\bigl(\cos(k_x+\Phi)+\cos(k_y-\Phi)\bigr).
	\end{equation}
	
	Hier spielen $v_A,v_B,v_C$ die Rolle effektiver Kanalgeschwindigkeiten.
	
	\subsection{Berry-Krümmung und Chern-Zahl}
	
	Für das untere Band wird numerisch nach Fukui--Hatsugai--Suzuki die Chern-Zahl berechnet. Formal gilt
	\begin{equation}
		C = \frac{1}{2\pi}\int_{\T^2} F_{12}(k)\,dk_x\,dk_y,
	\end{equation}
	mit der Berry-Krümmung $F_{12}$ des unteren Eigenraumbündels.
	
	Die globale von-Klitzing-artige Antwort wäre dann
	\begin{equation}
		\sigma_{xy}=\frac{e^2}{h}C.
	\end{equation}
	
	\subsection{Numerisches Bild}
	
	Der periodische Scan in den Parametern $(M,\Phi)$ zeigte diskrete Chern-Plateaus, insbesondere die Werte
	\begin{equation}
		C\in\{-1,0,1\}
	\end{equation}
	in den berücksichtigten Bereichen. Damit wird klar:
	\begin{itemize}
		\item Eulerzahl beschreibt die Topologie des diskreten Trägers,
		\item Aharonov--Bohm-Flüsse beschreiben Schleifenholonomie,
		\item Chern-Zahlen beschreiben die Topologie des Spektralbündels.
	\end{itemize}
	
	\section{Der optimale feinstrukturverschobene Fluss und seine Chern-Umgebung}
	
	Besonders interessant ist der Flusswert, der im endlichen Nullstellenmodell die beste Feinstruktur lieferte. Numerisch ergab sich ein effektiver optimaler Fluss in der Größenordnung
	\begin{equation}
		\frac{\Phi_{\mathrm{eff}}}{\pi}\approx 0.431261.
	\end{equation}
	
	Für diesen Fluss wurde das periodische Modell erneut in $M$ gescannt. Dabei zeigte sich:
	\begin{itemize}
		\item der betreffende Fluss liegt nicht in einer vollständig trivialen topologischen Umgebung,
		\item vielmehr treten in seiner Nachbarschaft nichttriviale Chern-Plateaus mit $C=\pm 1$ auf,
		\item zugleich bleibt ein großer Teil des untersuchten Bereichs bei $C=0$.
	\end{itemize}
	
	\begin{bemerkung}
		Dies spricht dafür, dass der numerisch beste Nullstellenfluss nicht tief in einem maximal robusten Plateau liegt, sondern eher in einer topologisch sensiblen Übergangszone, in der nichttriviale Phasen erreichbar sind.
	\end{bemerkung}
	
	\section{Interpretation}
	
	Die bisherige Kette lässt sich nun in kompakter Form zusammenfassen:
	\begin{enumerate}[label=(\arabic*)]
		\item Tetraederschichten organisieren die diskrete Hüllenstruktur.
		\item Die gewichtete Primzahruhr liefert eine kanalsensitive Uhrphase und vor allem eine numerisch wirksame Uhramplitude.
		\item Ein Gauge-Dirac-Operator auf dem E0/ABC-Kanalraum trägt die chirale Dynamik.
		\item Die Primzahruhr wirkt produktiv als Kantenfilter, nicht als diagonaler Energieterm.
		\item Echte Aharonov--Bohm-Flüsse auf elementaren Plaquetten verbessern die lokale Feinstruktur.
		\item Ein Feinstruktur-Shift $2\pi/137$ wird numerisch erst dann sichtbar, wenn er als Verschiebung des echten Plaquettenflusses eingeht.
		\item In einer periodisierten Version treten Berry-Krümmung und Chern-Zahlen auf.
	\end{enumerate}
	
	Die stärkste Lesart lautet daher:
	\begin{quote}
		Die Feinstrukturkonstante erscheint in diesem Modell nicht als bloße Zusatzphase eines einzelnen Kanals, sondern als mikroskopische Justierung einer echten Schleifenholonomie. Dadurch verbessert sie die lokale spektrale Feinstruktur, ohne das globale Wachstum wesentlich zu verändern.
	\end{quote}
	
	\section{Eulerzahl, Aharonov--Bohm und von Klitzing}
	
	Die drei zentralen topologischen Ebenen des Modells sind nun klar getrennt:
	\begin{itemize}
		\item \textbf{Eulerzahl} $\chi$: kombinatorische Topologie des diskreten Trägers,
		\item \textbf{Aharonov--Bohm-Holonomie}: lokale Phasenakkumulation auf echten Schleifen,
		\item \textbf{Chern-Zahl}: globale Topologie des periodischen Spektralbündels.
	\end{itemize}
	
	Diese Trennung ist konzeptionell entscheidend. Eulerzahl und Chern-Zahl sind keine konkurrierenden, sondern komplementäre Invarianten: Die eine beschreibt die Geometrie des Raums, die andere die Topologie des Spektrums auf diesem Raum.
	
	\section{Ausblick}
	
	Der vorliegende Rahmen legt mehrere nächste Schritte nahe:
	\begin{itemize}
		\item eine systematische Verfeinerung des endlichen Gauge-Dirac-Modells um echte Pentagon-Schleifen statt bloßer Plaquetten,
		\item die Untersuchung, ob der Flussbereich des besten Nullstellenfits mit Topologieübergängen im periodischen Modell noch präziser zusammenfällt,
		\item die Einbettung anisotroper Kanalgeschwindigkeiten $v_A,v_B,v_C$ in das endliche Modell selbst,
		\item und die Formulierung einer geschlossenen Theorie, in der Primzahruhr, Feinstruktur-Shift und topologische Bandantwort nicht nur numerisch, sondern auch konzeptionell als Einheit erscheinen.
	\end{itemize}
	
	Der wichtigste heuristische Satz des Artikels lautet:
	\begin{quote}
		Die lokale Feinstruktur der Nullstellen reagiert im Modell auf echte Schleifenholonomien, während die globale topologische Antwort in periodisierten Versionen durch Chern-Zahlen beschrieben wird; die Feinstrukturkonstante $1/137$ wirkt dabei produktiv als mikroskopische Verschiebung des Plaquettenflusses.
	\end{quote}
	
	\begin{thebibliography}{99}
		
		\bibitem{AharonovBohm}
		Y. Aharonov, D. Bohm,
		Significance of electromagnetic potentials in the quantum theory,
		\emph{Phys. Rev.} \textbf{115} (1959), 485--491.
		
		\bibitem{Klitzing}
		K. von Klitzing, G. Dorda, M. Pepper,
		New method for high-accuracy determination of the fine-structure constant based on quantized Hall resistance,
		\emph{Phys. Rev. Lett.} \textbf{45} (1980), 494--497.
		
		\bibitem{Thouless}
		D. J. Thouless, M. Kohmoto, M. P. Nightingale, M. den Nijs,
		Quantized Hall conductance in a two-dimensional periodic potential,
		\emph{Phys. Rev. Lett.} \textbf{49} (1982), 405--408.
		
		\bibitem{Fukui}
		T. Fukui, Y. Hatsugai, H. Suzuki,
		Chern numbers in discretized Brillouin zone: Efficient method of computing (spin) Hall conductances,
		\emph{J. Phys. Soc. Japan} \textbf{74} (2005), 1674--1677.
		
		\bibitem{Nakahara}
		M. Nakahara,
		\emph{Geometry, Topology and Physics},
		2nd ed., Taylor \& Francis, Boca Raton, 2003.
		
		\bibitem{Baez}
		J. C. Baez,
		The octonions,
		\emph{Bull. Amer. Math. Soc.} \textbf{39} (2002), 145--205.
		
	\end{thebibliography}
	
\end{document}
